\documentclass[12pt]{article}
\usepackage[margin=1in]{geometry}
\usepackage[T1]{fontenc}
\usepackage[utf8]{inputenc}
\usepackage{lmodern}
\usepackage{setspace}
\usepackage{hyperref}
\usepackage{booktabs}
\usepackage{longtable}
\usepackage{array}
\usepackage{graphicx}
\usepackage{enumitem}
\usepackage{xcolor}
\usepackage{float}
\usepackage{caption}
\usepackage{titlesec}
\usepackage{fancyhdr}
\usepackage{amsmath}

\hypersetup{colorlinks=true,linkcolor=black,citecolor=black,urlcolor=blue}
\setstretch{1.15}
\setlength{\parskip}{0.5em}
\setlength{\parindent}{0pt}
\titleformat{\section}{\large\bfseries}{\thesection}{0.75em}{}
\titleformat{\subsection}{\normalsize\bfseries}{\thesubsection}{0.75em}{}
\pagestyle{fancy}
\fancyhf{}
\fancyhead[L]{QNX Quantum Secure OTA Update}
\fancyhead[L]{Comp 4900E}
\fancyfoot[C]{\thepage}

\begin{document}

\begin{titlepage}
    \centering
    {\Large Carleton University \par}
    \vspace{1.5cm}
    {\LARGE \textbf{QNX Quantum Secure OTA Update} \par}
    \vspace{0.5cm}
    {\large COMP4900E, Real Time Operating Systems \par}
    \vspace{1.25cm}
    {\large Final Project Report\par}
    \vspace{1.5cm}
    \begin{tabular}{ll}
        Ben Whitten & \texttt{BenWhitten@cmail.carleton.ca} \\
        Jacob Bonner & \texttt{JacobBonner@cmail.carleton.ca} \\
        Caitlin Wardle & \texttt{CaitlinWardle@cmail.carleton.ca} \\
    \end{tabular}
    \vfill
    {\large \today \par}
\end{titlepage}
\begin{abstract}
    This project investigates whether a post-quantum digital signature scheme can be integrated into a QNX over-the-air (OTA) update workflow without making the update process impractical for embedded devices. Our target platform is a Raspberry Pi 4 running QNX OS 8.0, and our core comparison is between three update modes, unsigned OTA, classically signed OTA using RSA and Post Quantum Computing OTA using ML-DSA. The system is designed around an A/B partition update model in which a device checks a remote server for a new image, downloads an update bundle, verifies its signature, writes the new image to the inactive partition and reboots into the updated system if the validation succeeds. The report presents the need for secure OTA, the quantum risks facing current signing techniques, the benefits of ML-DSA, the implementation architecture and the comparisons across algorithms. Our evaluation focuses on signature generation time, signature verification time, total OTA update time, CPU usage, RAM usage, signature size, and storage overhead.
\end{abstract}
\tableofcontents
\newpage
\section{Introduction}
Embedded systems in automotive, medical, industrial and IoT environments usually rely on remote software updates to remain secure and functional over their deployment lifespan. Over the air (OTA) updates allow new software images or configuration changes to be delivered over a network, so physical access of the device is not required. OTA typically involves packaging an update on a server, transmitting it to a device, verifying its integrity and authenticity on the device, installing it and finally rebooting into the new version. For embedded systems OTA is beneficial because it reduces maintenance cost, minimizes system downtime, allows for rapid deployment and satisfies consumer need for at home maintenance. However, OTA has one critical problem, if an attacker can forger or tamper with an update, then they can gain persistent control of the device, a potentially disastrous scenario. This risk is exacerbated by the long lifespans of many embedded devices, since cryptographic assumptions used today may no longer be secure over the full life of the product. Most OTA sustems rely on classical public key encryption schemes, like RSA or elliptic-curve cryptography. But the long term security of these algorithms is threatened by large scale quantum computing, since new methods, like Shor's algorithm would break the hard mathematical problems these encryption schemes currently depend on. As such it is important to investigate whether post quantum signature mechanisms can be integrated into OTA validation, without unacceptable cost to the devices performance. This project aims to address that question by designing and evaluating a secure OTA mechanism for QNX based embedded systems using a Post Quantum Computing (PQC) signature algorithm for update validation. Our implementation is for QNX OS 8.0 on a Rasberry Pi 4, and uses ML-DSA through liboqs as the pqc signature scheme, with RSA through OpenSSL as the classical comparison baseline. As a baseline we have also included unsigned OTA, to isolate the cost of signature validation from the cost of the update process. 
\subsection{Research Question}
Is it feasible to replace a classical code signing algorithm in a QNX OTA update with a post quantum computing signature scheme without causing significant performance degradation on a simple embedded device?
\subsection{Objectives}
\begin{enumerate}[leftmargin=2em]
    \item Implement a working OTA update prototype on QNX 8.0 for Raspberry Pi 4.
    \item Integrate a digital signature verification step into the OTA validation path.
    \item Implement RSA signing and verification as a baseline.
    \item Implement ML-DSA signing and verification using liboqs as the post quantum solution
    \item Compare unsigned, RSA signed and ML-DSA signed updates using timing and resource usage.
    \item Analyze the resource demands of post quantum computing signatures on OTA updates in embedded devices to determine a baseline of resource and time requirements. As individual embedded systems have specific requirements the praticality of these drains will need to be determined on a per system basis.  
\end{enumerate}
\subsection{Scope}
The scope of this project is the OTA image packaging, download, signature validation, installation to an inactive partition and the reboot on an embedded device. The focus is on code signing and verification, not on transport encryption, key exchange or large scale deployments. Within this scope the project has three parts. Firstly it builds an OTA pipeline for QNX on Raspberry Pi 4 using an A/B partition model. Secondly it builds both a classical (RSA) and a post quantum computing (ML-DSA) signature into the update validation mechanism. Finally, it collects data pursuant to algorithm cost and device level measurements like runtime, CPU load, memory consumption and storage overhead.  
\subsection{Contributions}
\begin{itemize}[leftmargin=2em]
    \item \textbf{Ben Whitten:} \textcolor{red}{fill in}
    \item \textbf{Jacob Bonner:} \textcolor{red}{fill in}
    \item \textbf{Caitlin Wardle:} \textcolor{red}{fill in}
\end{itemize}
\section{Background}
\subsection{OTA Updates in Embedded Systems}
OTA updates allow a device to receive a new software image over a network connection
rather than through physical access. A typical OTA pipeline packages an image on
a server, transfers the image to the embedded device, validates the update on the device, installs the update, and reboots into the new version. For embedded systems,
OTA is especially useful for remote or critical devices. 
Despite the benefits, OTA updates are more vulnerable to installation failures and malicious attacks. If a device loses power, writes an incomplete
image, or boots an invalid update, it may become unusable or “bricked.” For this reason,
OTA systems rely on rollback capabilities,saved versions, and cryptographic
verification. A secure OTA design must protect against accidental
failure and against malicious changes. 
\subsection{Why Quantum Resistance Matters}
The signature verification step in current OTA systems uses classic public key
cryptography. The device stores a trusted public key and uses it to verify that the received
update was signed by the trusted private key holder. This model works well today, but its
long term security depends on mathematical formulas that future quantum computers may be able to break.
Shor’s algorithm would make integer factorization and elliptic curve discrete
logarithms solvable, destroying the security of RSA and ECC.
Since many embedded systems have long life spans, a device currently in use may still rely on these algorithms long after quantum attacks become effective. The \textit{harvest now, forge later}
threat model proposes attackers can collect update data and system images now and break them in the future, when quantum computers are developed. \textcolor{red}{Continue/make better}
\subsection{Post Quantum Signatures}
Post quantum computing (PQC) refers to cryptographic algorithms designed to remain secure even against attackers with quantum computers. For digital signatures the major categories of PQC are lattice based, hash based and multivariate constructions. This project focuses on ML-DSA, a lattice based digital signature algorithm  standardized by NIST and supported 
by liboqs.
ML-DSA was chosen because it offers a strong balance between security and practical performance for constrained systems. Compared to other algorithms like SLH-DSA, ML-DSA
generally provides smaller signatures and faster signing and verification, which are important when working with embedded devices operating with limited resources.
\subsection{Classical Signatures}
To determine the baseline overhead of current validation mechanisms this project has also implemented a classical scheme (RSA) that reflects present day OTA validations. The project uses RSA via OpenSSL as our baseline  baseline classical signature mechanism. RSA is well understood, widely supported, and easy to implement in a comparable signing and verification workflow. 
\section{System Design}
\subsection{Hardware}
The implementation was preformed with a Rasberry PI 4 running QNX OS 8.0. The Raspberry Pi 4 was chosen because it is affordable, well documented and practical for testing simple embedded devices. Using real hardware allows for device level measurements to realistically capture the performance limits and operational constrains of actual embedded systems.  
\subsection{A/B Partition OTA Model}
The OTA system is designed around an A/B partition architecture. In this model,
partition A contains the currently running system image, while partition B is inactive and can
be used as the destination for a downloaded update. Firstly the device will check the remote server for a new image. If a new image is available, the device downloads the update package. Then the device verifies the digital signature associated with the update. If the verification is successful, the new image is written to inactive partition B. The boot configuration is changed so the device boots from the updated partition B. The system reboots into the new image. This design has excellent rollback capabilities, if the verification fails, installation is aborted. If the update is written but becomes unstable, the system can fall back to the previously working partition. This makes the A/B model robust, and ideal for OTA. 
\subsection{Threat Model}
\textcolor{red}{What can go wrong/how can malicious users exploit the system and how are we solving it}
\subsection{ Cryptographic Workflow}
The Cryptographic workflow consists of two sides, the server side packaging step and the device side verification step.
\textbf{Server Side:}
\begin{enumerate}[leftmargin=2em]
\item Build or package the QNX image using \texttt{mkqnximage}
\item Compute a SHA 256 digest of the image
\item Sign the digest using RSA or ML-DSA
\item Publish the update image, signature and metadata to the server
\end{enumerate}
\textbf{Device Side:}
\begin{enumerate}[leftmargin=2em]
\item Download the update package from server.
\item Recompute the SHA 256 digest of the image.
\item Verify the signature using the \textcolor{red}{stored} public key.
\item Abandon installation if verification fails.
\item If verification succeeds, write the image to partition B and reboot.
\end{enumerate}
\section{Implementation}
\subsection{Software}
The main technologies used in this project are QNX Momentics IDE, \texttt{mkqnximage}, liboqs, OpenSSL, Git/GitHub, and Raspberry Pi 4 hardware. QNX Momentics is used as the main development environment. \texttt{mkqnximage} is used to generate QNX operating system images. OpenSSL provides classical RSA signing and verification, while liboqs provides the ML-DSA implementation for post quantum computing validation. Github \textcolor{red}{host the remote server} and acts as storage and version control for the project implementation. 
\subsection{OTA Update}
The unsigned OTA serves as a baseline measurement. In this mode, the device checks for an update, downloads the image, and writes it to the inactive partition without performing any kind of validation. This will establish a minimum time and resource cost of the OTA process itself. Without this baseline this project cannot accurately determine the cost of update retrieval and image download from the cost of the update validation mechanism 
\subsection{RSA OTA Validation}
The RSA method continues the baseline OTA mechanism, but adds a signing step on the server and a verification step on the device. On the server a SHA 256 digest of the update image is signed with an RSA private key through OpenSSL. On the device, the update image is hashed again and the signature is verified  with the RSA public key before installation proceeds. This models the classical public key verification currently used in many secure update systems. It provides a comparison point for evaluating the increased cost of post quantum signatures. 
\subsection{ML-DSA OTA Validation}
The post quantum method follows the same logic as RSA but replaces the signature scheme with ML-DSA using liboqs. On the server, the image digest is signed using an ML-DSA private key. On the device, the update digest is recomputed and verified using the corresponding ML-DSA public key. This implementation aims to keep workflow as consistent as possible while only swapping the signature scheme, in order to get a clean comparison of algorithmic cost. The main difference between RSA and ML-DSA is signature size and key size \textcolor{red}{include finding (if it increases verification time/storage,etc)}
\subsection{Metrics}
The following metrics are collected for each update mode:
\begin{itemize}[leftmargin=2em]
    \item \textbf{Signature Generation Time}: Time required to create the signature on the server side.
    \item \textbf{Signature Verification Time}: Time required to verify the signature on the device. 
    \item \textbf{Total OTA Time}: Full Time required from checking for updates/downloading update through validation to installation.
    \item \textbf{CPU Usage}: CPU cost observed during validation.
     \item \textbf{RAM Usage}: Maximum memory consumed during validation.
     \item \textbf{Signature Size}: Size in bytes of the generated signature.
     \item \textbf{Storage Overhead}: Additional package storage needed for signature components and keys. 
\end{itemize}
\subsection{Measurement Procedure}
\textcolor{red}{How we collected data}
\subsection{Data Analysis Procedure}
\textcolor{red}{How we summarized/display data (how we made graphs, what we measured for and how)}
\section{OTA Tutorial}
\textbf{\textcolor{red}{This is pulled from example final report repo, so we will need to augment}}
\subsection{Hardware \& Software Requirements}
\noindent
\textbf{Github Repo:} \url{https://github.com/BenjaminWhitten/COMP4900-Quantum-Secure-OTA} \\
\textbf{Host Server PC:} Windows 11 (with WSL Ubuntu) \\
\textbf{QNX SDP:} Version 7.1 and 8.0 (requires BlackBerry QNX license) \\
\textbf{Micro Controller:} Raspberry Pi 4

\noindent
\textbf{Required Cables/Accessories:}
\begin{itemize}
    \item USB to TTL serial cable (to access the Pi’s UART pins for console)
    \item Micro SD card (32 GB)
\end{itemize}
\subsection{Install QNX SDP and Tools}
\textcolor{red}{I fill in}
\subsection{Configure QNX}
\textcolor{red}{I fill in}
\subsection{Windows 11 SSH Setup}
\textcolor{red}{I fill in}
\subsection{Compile U-boot}
\textcolor{red}{I fill in}
\subsection{Prepare Rasberry Pi 4}
\textcolor{red}{I fill in}
\subsection{Performing the OTA Update}
\textcolor{red}{I fill in}
\section{Results and Analysis}
\subsection{Summary of Findings}
Table~\ref{tab:summary-results} is the primary benchmark summary table. Final values should be inserted after the full benchmark harness has been run.

\begin{table}[H]
\centering
\caption{Summary of OTA Results (\textcolor{red}{fill in later})}
\label{tab:summary-results}
\begin{tabular}{lccccccc}
\toprule
 & Sign Time & Verify Time & Total OTA & CPU & RAM & Sig Size & Storage Overhead \\
 & (ms) & (ms) & (ms) & (\%) & (KB) & (bytes) & (bytes) \\
\midrule
Unsigned & N/A & N/A & TBD & TBD & TBD & 0 & 0 \\
RSA & TBD & TBD & TBD & TBD & TBD & TBD & TBD \\
ML-DSA & TBD & TBD & TBD & TBD & TBD & TBD & TBD \\
\bottomrule
\end{tabular}
\end{table}
\subsection{Signature Generation Time}
\textcolor{red}{talk about quantitative results and analyze why + a nice graph}
\subsection{Signature Verification Time}
\textcolor{red}{talk about quantitative results and analyze why + a nice graph}
\subsection{Total OTA Time}
\textcolor{red}{talk about quantitative results and analyze why + a nice graph}
\subsection{Resource Usage}
\textcolor{red}{talk about quantitative results and analyze why + a nice graph of CPU and a nice graph of RAM}
\subsection{Signature Storage and Overhead}
\textcolor{red}{talk about quantitative results and analyze why + a nice graph}
\section{Discussion of Findings}
\subsection{Practical Meaning of the Results}
\subsection{Limitations}
\subsection{Future Work}
\section{Conclusion}
\Section{References}
\end{document}
